\documentclass[12pt]{article}
\usepackage[utf8]{inputenc}
\usepackage[T1]{fontenc}
\usepackage[spanish]{babel}
\shorthandoff{"}
\usepackage{csquotes}
\usepackage{amsmath,amsthm,bm}
\usepackage{eso-pic}
\usepackage{mathpazo}
\usepackage{setspace}
\setstretch{1.2}
\setlength{\parskip}{10pt}
\usepackage{indentfirst}
\usepackage{geometry}
\usepackage{graphicx}
\usepackage{titlesec}
\usepackage{fancyhdr}
\usepackage{ragged2e}
\usepackage{tikz}
\usetikzlibrary{arrows.meta,positioning}
\usepackage{subcaption}
\usepackage{mfirstuc}
\usepackage{appendix}
\usepackage{titletoc}
\usepackage{listings}
\usepackage{xcolor}
\usepackage{hyperref}
\hypersetup{hidelinks}
\usepackage[style=apa,backend=biber,sorting=none]{biblatex}
\addbibresource{export.bib}

\newtheorem{teorema}{Teorema}

\graphicspath{{./figuras/}}

\geometry{top=2.54cm, left=2.54cm, right=2.54cm, bottom=2.54cm}

\pagestyle{fancy}
\lhead{} \chead{} \rhead{\thepage}
\renewcommand{\headrulewidth}{0pt}
\lfoot{} \cfoot{} \rfoot{}
\renewcommand{\footrulewidth}{0pt}

\fancypagestyle{plain}{%
  \fancyhf{}%
  \renewcommand{\headrulewidth}{0pt}%
  \renewcommand{\footrulewidth}{0pt}%
}

% Configuración de estilo para el código MATLAB
\lstset{
    language=Matlab,
    basicstyle=\ttfamily\small,
    keywordstyle=\color{blue}\bfseries,
    commentstyle=\color{green!60!black},
    stringstyle=\color{purple},
    numbers=left,
    numberstyle=\tiny\color{gray},
    breaklines=true,
    frame=single,
    captionpos=b,
    showstringspaces=false,
    extendedchars=true,
    literate=%
    {á}{{\'a}}1 {é}{{\'e}}1 {í}{{\'i}}1 {ó}{{\'o}}1 {ú}{{\'u}}1
    {Á}{{\'A}}1 {É}{{\'E}}1 {Í}{{\'I}}1 {Ó}{{\'O}}1 {Ú}{{\'U}}1
    {ñ}{{\~n}}1 {Ñ}{{\~N}}1
}

\begin{document}

\begin{titlepage}
    \centering
    \singlespacing
    \vspace*{\fill}
    \vspace{-2cm}
    \includegraphics[width=3.5cm]{logo1.png}\par
    \vspace{0.1cm}
    {\large Universidad Distrital Francisco José de Caldas \par}
    \vspace{0.2cm}
    {\large Facultad de Ingeniería \\Ingeniería Catastral y Geodesia\par}
    \vspace{1.5cm}
    {\large \bfseries
    Tarea 10 \\Herramientas Matemáticas \\para el Procesamiento de la Información \par}
    \vspace{1.5cm}
    {\large Presentado por \\Argenis Alexandra Daza Roa \par}
    \vspace{1.5cm}
    {\large Docente \\Msc. Hans Igor López Chávez \par}
    \vspace{1.5cm}
    {\large 21 de octubre de 2026 \par}
    \vspace{0.5cm}
    {\large Bogotá D.C. \par}
    \vspace*{\fill}
\end{titlepage}

\newpage

\section{Desarrollo}

Individualmente, desarrolle el modelo matemático de la cola M/M/1/K. En concreto:

\begin{itemize}
    \item Plantee el diagrama de transición de estados.
    \item Obtenga las ecuaciones de estado estacionario.
    \item Obtenga $P_n$ en función de $P_0$.
    \item Obtenga $P_n$ en función de los parámetros $\lambda$ y $\mu$.
\end{itemize}

\subsection{Descripción del sistema}

Según la notación de Kendall-Lee $A/B/C/D/E/F$ \autocite[cap.~1]{shortle2018}, la cola M/M/1/K corresponde a M/M/1/K/$\infty$/FIFO (población infinita, disciplina FIFO), es decir:

\begin{itemize}
    \item Llegadas de Poisson con tasa $\lambda$ (tiempos entre llegadas $T\sim \mathrm{Exp}(x;\lambda)$).
    \item Un servidor con tiempos de servicio exponenciales de tasa $\mu$ (sin memoria, $T\sim \mathrm{Exp}(x;\mu)$).
    \item Capacidad máxima del sistema igual a $K$ (línea de espera más servidor): un cliente que llega cuando hay $K$ en el sistema es rechazado.
\end{itemize}

Sea $P_n$ la probabilidad de que haya $n$ clientes en el sistema, con $n=0,1,\dots,K$.

\begin{figure}[h]
\centering
\begin{tikzpicture}[>=Stealth, every node/.style={font=\small}]
    \draw (1.2,0) rectangle (4.2,0.8);
    \draw (2.2,0) -- (2.2,0.8);
    \draw (3.2,0) -- (3.2,0.8);
    \node at (1.7,0.4) {$\cdots$};
    \draw (5.0,0.4) circle (0.45);
    \node at (5.0,0.4) {$\mu$};
    \draw[->] (4.2,0.4) -- (4.55,0.4);
    \draw[->] (5.45,0.4) -- (6.3,0.4);
    \draw[->] (-0.6,0.4) -- node[above]{$\lambda$} (1.2,0.4);
    \draw[<->] (1.2,-0.35) -- node[below]{$K-1$ espacios} (4.2,-0.35);
\end{tikzpicture}
\caption{Esquema del sistema M/M/1/K.}
\end{figure}

\subsection{Diagrama de transición de estados}

El sistema es un proceso de nacimiento y muerte con espacio de estados finito $\{0,1,\dots,K\}$ \autocite[cap.~3]{shortle2018}. Un nacimiento ($n\to n+1$) ocurre con tasa $\lambda$ solo si $n<K$, pues en el estado $K$ el sistema está lleno y las llegadas se pierden. Una muerte ($n\to n-1$) ocurre con tasa $\mu$ si $n\geq 1$. La tasa promedio de cada transición es la tasa multiplicada por la probabilidad del estado de origen.

\begin{figure}[h]
\centering
\begin{tikzpicture}[>=Stealth, node distance=2.15cm,
    st/.style={circle,draw,minimum size=0.95cm}]
    \node[st] (0) {0};
    \node[st,right=of 0] (1) {1};
    \node[st,right=of 1] (2) {2};
    \node[right=of 2] (d) {$\cdots$};
    \node[st,right=of d] (k1) {\scriptsize$K{-}1$};
    \node[st,right=of k1] (k) {$K$};

    \draw[->,blue] (0) to[bend left=40] node[above]{$\lambda P_0$} (1);
    \draw[->,blue] (1) to[bend left=40] node[above]{$\lambda P_1$} (2);
    \draw[->,blue] (2) to[bend left=40] node[above]{$\lambda P_2$} (d);
    \draw[->,blue] (d) to[bend left=40] node[above]{\scriptsize$\lambda P_{K-2}$} (k1);
    \draw[->,blue] (k1) to[bend left=40] node[above]{\scriptsize$\lambda P_{K-1}$} (k);

    \draw[->,red] (1) to[bend left=40] node[below]{$\mu P_1$} (0);
    \draw[->,red] (2) to[bend left=40] node[below]{$\mu P_2$} (1);
    \draw[->,red] (d) to[bend left=40] node[below]{$\mu P_3$} (2);
    \draw[->,red] (k1) to[bend left=40] node[below]{\scriptsize$\mu P_{K-1}$} (d);
    \draw[->,red] (k) to[bend left=40] node[below]{$\mu P_K$} (k1);
\end{tikzpicture}
\caption{Diagrama de transición de estados de la cola M/M/1/K.}
\end{figure}

\subsection{Ecuaciones de estado estacionario}

En estado estacionario, para cada estado la tasa promedio de salida del estado es igual a la tasa promedio de llegada al estado (ecuaciones de balance) \autocite[cap.~3]{shortle2018}.

\begin{center}
\begin{tabular}{c c c c}
\hline
Estado & Tasa promedio de salida del estado & $=$ & Tasa promedio de llegada al estado \\
\hline
$0$ & $\lambda P_0$ & $=$ & $\mu P_1$ \\[4pt]
$1$ & $\mu P_1+\lambda P_1$ & $=$ & $\lambda P_0+\mu P_2$ \\[4pt]
$2$ & $\mu P_2+\lambda P_2$ & $=$ & $\lambda P_1+\mu P_3$ \\[4pt]
$\vdots$ & $\vdots$ & & $\vdots$ \\[4pt]
$n$ & $\mu P_n+\lambda P_n$ & $=$ & $\lambda P_{n-1}+\mu P_{n+1}$ \quad $(1\le n\le K-1)$ \\[4pt]
$K$ & $\mu P_K$ & $=$ & $\lambda P_{K-1}$ \\
\hline
\end{tabular}
\end{center}

En el estado $K$ no aparece el término de salida $\lambda P_K$, porque desde $K$ no hay nacimientos (el sistema está lleno).

\subsection{$P_n$ en función de $P_0$}

Del estado $0$:
\[
\lambda P_0=\mu P_1 \;\Rightarrow\; P_1=\frac{\lambda}{\mu}P_0 .
\]

Del estado $1$:
\[
\mu P_1+\lambda P_1=\lambda P_0+\mu P_2 .
\]
Como $\lambda P_0=\mu P_1$, se cancelan $\mu P_1$ y $\lambda P_0$ y queda
\[
\lambda P_1=\mu P_2 \;\Rightarrow\; P_2=\frac{\lambda}{\mu}P_1=\left(\frac{\lambda}{\mu}\right)^2P_0 .
\]

Del estado $2$, usando $\lambda P_1=\mu P_2$:
\[
\lambda P_2=\mu P_3 \;\Rightarrow\; P_3=\frac{\lambda}{\mu}P_2=\left(\frac{\lambda}{\mu}\right)^3P_0 .
\]

En general, cada ecuación, una vez simplificada con la anterior, da:
\[
\lambda P_n=\mu P_{n+1}\;\Rightarrow\; P_{n+1}=\frac{\lambda}{\mu}P_n,\qquad n=0,1,\dots,K-1 .
\]
Es decir, la probabilidad de un estado depende solo del estado anterior. Por inducción: si $P_n=\left(\frac{\lambda}{\mu}\right)^nP_0$, entonces $P_{n+1}=\frac{\lambda}{\mu}P_n=\left(\frac{\lambda}{\mu}\right)^{n+1}P_0$. Por lo tanto,
\[
\boxed{P_n=\left(\frac{\lambda}{\mu}\right)^{n}P_0,\qquad n=0,1,\dots,K}
\]
La ecuación del estado $K$, $\mu P_K=\lambda P_{K-1}$, coincide con la recurrencia anterior para $n=K-1$; es redundante, pero confirma que la fórmula también es válida en el último estado $n=K$.

\subsection{$P_n$ en función de $\lambda$ y $\mu$}

Sea $A=\dfrac{\lambda}{\mu}$ (intensidad de tráfico). Como el sistema solo puede estar en los estados $0,\dots,K$, la condición de normalización es
\[
\sum_{n=0}^{K}P_n=1
\;\Rightarrow\;
P_0\sum_{n=0}^{K}A^n=1
\;\Rightarrow\;
P_0=\frac{1}{S_K},\qquad S_K=\sum_{n=0}^{K}A^n ,
\]
donde $S_K$ es la suma de los primeros $K+1$ términos.

\textbf{Cálculo de $S_K$} (suma geométrica finita):
\begin{align*}
S_K &= 1+A+A^2+\cdots+A^{K-1}+A^K \\
AS_K &= \phantom{1+{}}A+A^2+\cdots+A^{K-1}+A^K+A^{K+1}
\end{align*}
Restando ambas expresiones se cancelan todos los términos intermedios:
\[
S_K-AS_K=1-A^{K+1}\;\Rightarrow\;S_K=\frac{1-A^{K+1}}{1-A},\qquad A\neq 1 .
\]

Entonces
\[
\boxed{P_0=\frac{1-A}{1-A^{K+1}}=\frac{1-\frac{\lambda}{\mu}}{1-\left(\frac{\lambda}{\mu}\right)^{K+1}}}
\]
y sustituyendo en $P_n=A^nP_0$:
\[
\boxed{P_n=\frac{\left(1-\frac{\lambda}{\mu}\right)\left(\frac{\lambda}{\mu}\right)^{n}}{1-\left(\frac{\lambda}{\mu}\right)^{K+1}}=\frac{(1-A)A^n}{1-A^{K+1}},\qquad n=0,1,\dots,K,\quad \lambda\neq\mu}
\]

De forma equivalente, multiplicando numerador y denominador por $\mu^{K+1}$:
\[
P_n=\frac{(\mu-\lambda)\,\lambda^{n}\,\mu^{K-n}}{\mu^{K+1}-\lambda^{K+1}},\qquad n=0,1,\dots,K .
\]

Este resultado coincide con el caso de un servidor ($c=1$) de la cola de capacidad finita M/M/c/K \autocite[cap.~3]{shortle2018}.

\textbf{Relación con la M/M/1.} En la M/M/1 (capacidad infinita) la suma es infinita, $S_\infty=\frac{1}{1-A}$, la cual converge solo si $|A|<1$, es decir $\lambda<\mu$; así se obtiene $P_0=1-A$ y $P_n=\left(\frac{\lambda}{\mu}\right)^n\left(1-\frac{\lambda}{\mu}\right)$. En la M/M/1/K la suma es finita y por eso aparece el denominador $1-A^{K+1}$, que no se puede eliminar.

\textbf{Caso $\lambda=\mu$ ($A=1$).} Entonces $S_K=\sum_{n=0}^{K}1=K+1$, de modo que
\[
P_0=P_n=\frac{1}{K+1},\qquad n=0,1,\dots,K .
\]

\textbf{Observaciones.}
\begin{itemize}
    \item Como la suma es finita, el sistema alcanza estado estacionario para cualquier $\lambda,\mu>0$; no se requiere $\lambda<\mu$ (condición que sí exige la M/M/1 para que la serie converja, $|A|<1$).
    \item Si $\lambda<\mu$ y $K\to\infty$, entonces $A^{K+1}\to 0$ y se recupera la M/M/1: $P_n=\left(\frac{\lambda}{\mu}\right)^n\left(1-\frac{\lambda}{\mu}\right)$.
    \item La probabilidad de bloqueo (cliente rechazado) es $P_K=\dfrac{(1-A)A^K}{1-A^{K+1}}$.
\end{itemize}

\printbibliography

\end{document}
